\documentclass[11pt,a4paper]{article}
\usepackage[utf8]{inputenc}
\usepackage[T1]{fontenc}
\usepackage[margin=2.2cm]{geometry}
\usepackage{amsmath,amssymb}
\usepackage{enumitem}
\usepackage{hyperref}

\setlength{\parindent}{0pt}
\setlength{\parskip}{0.75em}

\title{MA087G Mathematical Statistics\\Beginner-Friendly Full Solution Guide\\2026-01-07}
\author{}
\date{}

\begin{document}
\maketitle

\section*{Source notes before solving}
The official exam PDF was image-only. That means normal text extraction found no text, but the rendered page images were readable.

Two details are important:
\begin{itemize}[leftmargin=*]
    \item Uppgift 2 has an internal mismatch. The main displayed density is \(f(x)=a x^2(2-x)\), but the printed fallback comment says \(f(x)=5x^3(2-x)/8\). I solve the main displayed density, because that is the actual formula in the question. I also give the fallback answers at the end of Uppgift 2.
    \item Uppgift 4b does not explicitly say that \(X\) and \(Y\) are independent. To calculate the variance of the taxi cost, we need independence. I use the standard exam assumption that they are independent.
\end{itemize}

\section*{Uppgift 1}

\subsection*{What this question is really asking}
We have two events, \(A\) and \(B\). An event just means something that may or may not happen.

The question gives three pieces of information:
\[
P(A^c\cap B^c)=0.60,\qquad P(A\cap B)=0.20,\qquad P(A^c\cap B)=0.10.
\]

The little \(c\) means ``not''. So:
\[
A^c=\text{not }A,\qquad B^c=\text{not }B.
\]

The question asks:
\begin{itemize}[leftmargin=*]
    \item Are \(A\) and \(B\) independent?
    \item What is \(P(B\mid A^c)\), the probability that \(B\) happens given that \(A\) does not happen?
\end{itemize}

\subsection*{How to recognize this type of question}
The clues are:
\begin{itemize}[leftmargin=*]
    \item It talks about two events \(A\) and \(B\).
    \item It gives probabilities for combinations like \(A^c\cap B\).
    \item It asks about independence.
    \item It asks for a conditional probability, shown by the vertical bar \(\mid\).
\end{itemize}

So this is a two-event probability table problem. The safest beginner method is to fill in the four possible boxes:
\[
A\cap B,\quad A\cap B^c,\quad A^c\cap B,\quad A^c\cap B^c.
\]

\subsection*{Tiny theory needed}
The symbol \(\cap\) means ``and''. For example:
\[
A\cap B=\text{\(A\) happens and \(B\) happens}.
\]

The symbol \(P(B\mid A^c)\) means:
\[
\text{probability of \(B\), if we already know that \(A\) did not happen}.
\]

The formula for conditional probability is:
\[
P(B\mid A^c)=\frac{P(B\cap A^c)}{P(A^c)}.
\]

Independence means that knowing one event happened does not change the probability of the other. The exam formula for checking independence is:
\[
P(A\cap B)=P(A)P(B).
\]

\subsection*{Formula and what everything means}
For independence:
\[
P(A\cap B)=P(A)P(B).
\]
Here:
\begin{itemize}[leftmargin=*]
    \item \(P(A\cap B)\) is the probability that both happen.
    \item \(P(A)\) is the probability that \(A\) happens.
    \item \(P(B)\) is the probability that \(B\) happens.
    \item If left side and right side are equal, they are independent.
    \item If they are not equal, they are not independent.
\end{itemize}

For conditional probability:
\[
P(B\mid A^c)=\frac{P(B\cap A^c)}{P(A^c)}.
\]
Here:
\begin{itemize}[leftmargin=*]
    \item The top is the part where \(B\) happens and \(A\) does not happen.
    \item The bottom is all cases where \(A\) does not happen.
\end{itemize}

\subsection*{Step-by-step solution}
The four possible boxes must add to \(1\), because one of them must happen:
\[
P(A\cap B)+P(A\cap B^c)+P(A^c\cap B)+P(A^c\cap B^c)=1.
\]

We already know three boxes:
\[
P(A\cap B)=0.20,
\]
\[
P(A^c\cap B)=0.10,
\]
\[
P(A^c\cap B^c)=0.60.
\]

Find the missing box:
\begin{align*}
P(A\cap B^c)
&=1-0.20-0.10-0.60\\
&=1-0.90\\
&=0.10.
\end{align*}
This means the probability that \(A\) happens but \(B\) does not happen is \(0.10\).

Now find \(P(A)\). Event \(A\) happens in two boxes:
\[
A\cap B
\]
and
\[
A\cap B^c.
\]
So:
\begin{align*}
P(A)
&=P(A\cap B)+P(A\cap B^c)\\
&=0.20+0.10\\
&=0.30.
\end{align*}

Now find \(P(B)\). Event \(B\) happens in two boxes:
\[
A\cap B
\]
and
\[
A^c\cap B.
\]
So:
\begin{align*}
P(B)
&=P(A\cap B)+P(A^c\cap B)\\
&=0.20+0.10\\
&=0.30.
\end{align*}

Now check independence:
\begin{align*}
P(A)P(B)
&=0.30\cdot 0.30\\
&=0.09.
\end{align*}
But:
\[
P(A\cap B)=0.20.
\]
Since:
\[
0.20\neq 0.09,
\]
\(A\) and \(B\) are not independent.

Now calculate \(P(B\mid A^c)\):
\begin{align*}
P(B\mid A^c)
&=\frac{P(B\cap A^c)}{P(A^c)}.
\end{align*}
The top is given:
\[
P(B\cap A^c)=P(A^c\cap B)=0.10.
\]
The bottom \(P(A^c)\) means all cases where \(A\) does not happen:
\begin{align*}
P(A^c)
&=P(A^c\cap B)+P(A^c\cap B^c)\\
&=0.10+0.60\\
&=0.70.
\end{align*}
Therefore:
\begin{align*}
P(B\mid A^c)
&=\frac{0.10}{0.70}\\
&=\frac{1}{7}\\
&\approx 0.143.
\end{align*}

\subsection*{Answer}
\[
\boxed{A \text{ and } B \text{ are not independent}}
\]
\[
\boxed{P(B\mid A^c)=\frac{1}{7}\approx 0.143}
\]

\subsection*{What to memorize for the exam}
For independence, always check:
\[
P(A\cap B)=P(A)P(B).
\]
For conditional probability, always use:
\[
P(\text{wanted}\mid \text{given})=
\frac{P(\text{wanted and given})}{P(\text{given})}.
\]

\subsection*{Common trap}
Do not compare \(P(A\cap B)\) with \(P(A)+P(B)\). Independence uses multiplication, not addition.

\subsection*{Repeated exam pattern}
This is the same broad pattern as 2025-08, Uppgift 1: two-event probability with complements and conditional probability. The numbers changed, but the solving method is the same.

\section*{Uppgift 2}

\subsection*{What this question is really asking}
We are given a function \(f(x)\) and need to make it into a probability density function.

A density function is used for a continuous random variable. Continuous means \(X\) can take values on an interval, like between \(0\) and \(2\), instead of only separate values like \(0,1,2\).

The question asks us to:
\begin{itemize}[leftmargin=*]
    \item find the constant \(a\),
    \item find the distribution function \(F_X(x)\),
    \item calculate \(P(X>1)\),
    \item calculate expected value and standard deviation.
\end{itemize}

\subsection*{How to recognize this type of question}
The clues are:
\begin{itemize}[leftmargin=*]
    \item The question says ``täthetsfunktion'', which means density function.
    \item The function is given on an interval \(0\le x\le 2\).
    \item It asks for \(F_X\), which is the distribution function/CDF.
    \item It asks for expected value and standard deviation.
\end{itemize}

So this is an integral problem. For continuous random variables, probabilities are areas, and areas are found with integrals.

\subsection*{Tiny theory needed}
The displayed density is:
\[
f(x)=
\begin{cases}
a x^2(2-x), & 0\le x\le 2,\\
0, & \text{otherwise}.
\end{cases}
\]

For \(f\) to be a density, the total area under the curve must be \(1\):
\[
\int_{-\infty}^{\infty}f(x)\,dx=1.
\]
Since \(f(x)=0\) outside \(0\le x\le 2\), we only integrate from \(0\) to \(2\).

The distribution function \(F_X(x)\) means:
\[
F_X(x)=P(X\le x).
\]
It tells us how much probability is to the left of \(x\).

\subsection*{Formula and what everything means}
Density condition:
\[
\int_0^2 f(x)\,dx=1.
\]
Here:
\begin{itemize}[leftmargin=*]
    \item \(f(x)\) is the density.
    \item The integral is the area under \(f(x)\).
    \item The total area must equal \(1\), because total probability is \(100\%\).
\end{itemize}

CDF formula:
\[
F_X(x)=\int_{-\infty}^{x}f(t)\,dt.
\]
Here:
\begin{itemize}[leftmargin=*]
    \item \(F_X(x)\) is probability up to \(x\).
    \item \(t\) is only a temporary integration letter.
\end{itemize}

Expected value:
\[
E(X)=\int x f(x)\,dx.
\]
This is the long-run average value.

Variance:
\[
V(X)=E(X^2)-[E(X)]^2.
\]
Standard deviation:
\[
D(X)=\sqrt{V(X)}.
\]

\subsection*{Step-by-step solution}
\textbf{Part a: Find \(a\).}

Use the density condition:
\begin{align*}
\int_0^2 a x^2(2-x)\,dx&=1.
\end{align*}
First expand inside the integral:
\[
x^2(2-x)=2x^2-x^3.
\]
So:
\begin{align*}
a\int_0^2(2x^2-x^3)\,dx&=1.
\end{align*}
Now integrate:
\[
\int 2x^2\,dx=\frac{2x^3}{3},
\qquad
\int x^3\,dx=\frac{x^4}{4}.
\]
So:
\begin{align*}
a\left[\frac{2x^3}{3}-\frac{x^4}{4}\right]_0^2&=1.
\end{align*}
Insert the upper limit \(2\):
\begin{align*}
\frac{2\cdot 2^3}{3}-\frac{2^4}{4}
&=\frac{16}{3}-4\\
&=\frac{16}{3}-\frac{12}{3}\\
&=\frac{4}{3}.
\end{align*}
Insert the lower limit \(0\), which gives \(0\). Therefore:
\begin{align*}
a\cdot \frac{4}{3}&=1\\
a&=\frac{3}{4}.
\end{align*}

\textbf{Part b: Find \(F_X(x)\).}

The CDF must be piecewise:
\begin{itemize}[leftmargin=*]
    \item If \(x<0\), we are before the density starts, so \(F_X(x)=0\).
    \item If \(0\le x\le 2\), integrate from \(0\) to \(x\).
    \item If \(x>2\), we have passed the whole density, so \(F_X(x)=1\).
\end{itemize}

For \(0\le x\le 2\):
\begin{align*}
F_X(x)
&=\int_0^x \frac{3}{4}t^2(2-t)\,dt\\
&=\frac{3}{4}\int_0^x(2t^2-t^3)\,dt\\
&=\frac{3}{4}\left[\frac{2t^3}{3}-\frac{t^4}{4}\right]_0^x\\
&=\frac{3}{4}\left(\frac{2x^3}{3}-\frac{x^4}{4}\right).
\end{align*}
Now multiply in \(\frac{3}{4}\):
\begin{align*}
F_X(x)
&=\frac{x^3}{2}-\frac{3x^4}{16}.
\end{align*}

Therefore:
\[
F_X(x)=
\begin{cases}
0, & x<0,\\[1mm]
\dfrac{x^3}{2}-\dfrac{3x^4}{16}, & 0\le x\le 2,\\[2mm]
1, & x>2.
\end{cases}
\]

\textbf{Part c: Calculate \(P(X>1)\).}

The CDF gives \(P(X\le 1)\). We want the opposite side:
\[
P(X>1)=1-P(X\le 1)=1-F_X(1).
\]
Calculate \(F_X(1)\):
\begin{align*}
F_X(1)
&=\frac{1^3}{2}-\frac{3\cdot 1^4}{16}\\
&=\frac{1}{2}-\frac{3}{16}\\
&=\frac{8}{16}-\frac{3}{16}\\
&=\frac{5}{16}.
\end{align*}
So:
\begin{align*}
P(X>1)
&=1-\frac{5}{16}\\
&=\frac{11}{16}.
\end{align*}

\textbf{Part d: Expected value and standard deviation.}

First calculate \(E(X)\):
\begin{align*}
E(X)
&=\int_0^2 x f(x)\,dx\\
&=\int_0^2 x\cdot \frac{3}{4}x^2(2-x)\,dx\\
&=\frac{3}{4}\int_0^2 x^3(2-x)\,dx\\
&=\frac{3}{4}\int_0^2(2x^3-x^4)\,dx.
\end{align*}
Now integrate:
\begin{align*}
E(X)
&=\frac{3}{4}\left[\frac{x^4}{2}-\frac{x^5}{5}\right]_0^2\\
&=\frac{3}{4}\left(\frac{2^4}{2}-\frac{2^5}{5}\right)\\
&=\frac{3}{4}\left(8-\frac{32}{5}\right)\\
&=\frac{3}{4}\cdot \frac{8}{5}\\
&=\frac{6}{5}.
\end{align*}

Now calculate \(E(X^2)\), because variance needs it:
\begin{align*}
E(X^2)
&=\int_0^2 x^2 f(x)\,dx\\
&=\int_0^2 x^2\cdot \frac{3}{4}x^2(2-x)\,dx\\
&=\frac{3}{4}\int_0^2(2x^4-x^5)\,dx\\
&=\frac{3}{4}\left[\frac{2x^5}{5}-\frac{x^6}{6}\right]_0^2\\
&=\frac{3}{4}\left(\frac{64}{5}-\frac{64}{6}\right)\\
&=\frac{3}{4}\left(\frac{64}{5}-\frac{32}{3}\right)\\
&=\frac{8}{5}.
\end{align*}

Now variance:
\begin{align*}
V(X)
&=E(X^2)-[E(X)]^2\\
&=\frac{8}{5}-\left(\frac{6}{5}\right)^2\\
&=\frac{8}{5}-\frac{36}{25}\\
&=\frac{40}{25}-\frac{36}{25}\\
&=\frac{4}{25}.
\end{align*}
Standard deviation:
\begin{align*}
D(X)
&=\sqrt{V(X)}\\
&=\sqrt{\frac{4}{25}}\\
&=\frac{2}{5}.
\end{align*}

\subsection*{Answer}
For the displayed density \(a x^2(2-x)\):
\[
\boxed{a=\frac{3}{4}}
\]
\[
\boxed{
F_X(x)=
\begin{cases}
0, & x<0,\\
\frac{x^3}{2}-\frac{3x^4}{16}, & 0\le x\le 2,\\
1, & x>2
\end{cases}}
\]
\[
\boxed{P(X>1)=\frac{11}{16}}
\]
\[
\boxed{E(X)=\frac{6}{5}},\qquad
\boxed{D(X)=\frac{2}{5}}.
\]

If the printed fallback density \(5x^3(2-x)/8\) was intended instead, then:
\[
F_X(x)=\frac{5x^4}{16}-\frac{x^5}{8},\qquad
P(X>1)=\frac{13}{16},
\]
\[
E(X)=\frac{4}{3},\qquad
D(X)=\sqrt{\frac{8}{63}}.
\]

\subsection*{What to memorize for the exam}
For density questions:
\begin{itemize}[leftmargin=*]
    \item First make total area equal to \(1\).
    \item Then find the CDF by integrating the density.
    \item For \(P(X>c)\), use \(1-F(c)\).
    \item For standard deviation, calculate variance first.
\end{itemize}

\subsection*{Common trap}
Do not write only the middle part of the CDF. A CDF must also say what happens before and after the interval.

\subsection*{Repeated exam pattern}
This is the same pattern as 2025-08, Uppgift 2: continuous density normalization, CDF, probability, expected value, and standard deviation. The density changed, but the solving recipe is identical.

\section*{Uppgift 3}

\subsection*{What this question is really asking}
We have two independent exponential random variables:
\[
X\sim \operatorname{Exp}(0.04),\qquad Y\sim \operatorname{Exp}(0.05).
\]

The new variable is:
\[
Z=\min(X,Y).
\]
That means \(Z\) is whichever of \(X\) and \(Y\) happens first/smaller.

The question asks:
\begin{itemize}[leftmargin=*]
    \item What distribution does \(Z\) have?
    \item What is \(P(Z\ge 10)\)?
    \item What is \(E(Z)\)?
\end{itemize}

\subsection*{How to recognize this type of question}
The clues are:
\begin{itemize}[leftmargin=*]
    \item It says \(X\in \operatorname{Exp}(0.04)\) and \(Y\in \operatorname{Exp}(0.05)\).
    \item It says independent, which means we may multiply probabilities.
    \item It defines \(Z=\min(X,Y)\).
\end{itemize}

For a minimum, use survival probabilities. Survival means probability of being larger than a value.

\subsection*{Tiny theory needed}
If \(X\sim \operatorname{Exp}(\lambda)\), then:
\[
P(X>t)=e^{-\lambda t}.
\]
Here \(\lambda\) is the rate.

For \(Z=\min(X,Y)\), the event \(Z>t\) means both variables are bigger than \(t\):
\[
Z>t
\]
means:
\[
X>t \text{ and } Y>t.
\]

Because \(X\) and \(Y\) are independent, we can multiply:
\[
P(X>t \text{ and } Y>t)=P(X>t)P(Y>t).
\]

\subsection*{Formula and what everything means}
Exponential survival formula:
\[
P(X>t)=e^{-\lambda t}.
\]
Here:
\begin{itemize}[leftmargin=*]
    \item \(t\) is the time or value we compare against.
    \item \(\lambda\) is the exponential rate.
    \item \(e^{-\lambda t}\) gives the probability that the variable is still above \(t\).
\end{itemize}

For a minimum of independent exponentials:
\[
\min(\operatorname{Exp}(\lambda_1),\operatorname{Exp}(\lambda_2))
\sim \operatorname{Exp}(\lambda_1+\lambda_2).
\]
In words: for the minimum, the rates add.

\subsection*{Step-by-step solution}
\textbf{Part a: Find the distribution of \(Z\).}

Start with:
\[
Z=\min(X,Y).
\]
For \(z\ge 0\):
\begin{align*}
P(Z>z)
&=P(\min(X,Y)>z).
\end{align*}
The minimum is greater than \(z\) only if both \(X\) and \(Y\) are greater than \(z\):
\begin{align*}
P(Z>z)
&=P(X>z,\;Y>z).
\end{align*}
Since \(X\) and \(Y\) are independent:
\begin{align*}
P(Z>z)
&=P(X>z)P(Y>z).
\end{align*}
Use the exponential survival formula:
\begin{align*}
P(X>z)&=e^{-0.04z},\\
P(Y>z)&=e^{-0.05z}.
\end{align*}
Multiply:
\begin{align*}
P(Z>z)
&=e^{-0.04z}e^{-0.05z}\\
&=e^{-0.09z}.
\end{align*}
This is the survival function of an exponential variable with rate \(0.09\). Therefore:
\[
Z\sim \operatorname{Exp}(0.09).
\]

If we want the distribution function:
\[
F_Z(z)=P(Z\le z).
\]
Since:
\[
P(Z\le z)=1-P(Z>z),
\]
we get:
\[
F_Z(z)=
\begin{cases}
0, & z<0,\\
1-e^{-0.09z}, & z\ge 0.
\end{cases}
\]

\textbf{Part b: Calculate \(P(Z\ge 10)\).}

For a continuous variable, \(P(Z\ge 10)\) and \(P(Z>10)\) are the same, because the probability of exactly \(10\) is \(0\).

Use the survival function:
\begin{align*}
P(Z\ge 10)
&=P(Z>10)\\
&=e^{-0.09\cdot 10}\\
&=e^{-0.9}\\
&\approx 0.407.
\end{align*}

\textbf{Part c: Find \(E(Z)\).}

For an exponential variable:
\[
E(Z)=\frac{1}{\lambda}.
\]
Here \(\lambda=0.09\), so:
\begin{align*}
E(Z)
&=\frac{1}{0.09}\\
&=\frac{100}{9}\\
&\approx 11.11.
\end{align*}

\subsection*{Answer}
\[
\boxed{Z\sim \operatorname{Exp}(0.09)}
\]
\[
\boxed{P(Z\ge 10)=e^{-0.9}\approx 0.407}
\]
\[
\boxed{E(Z)=\frac{100}{9}\approx 11.11}
\]

\subsection*{What to memorize for the exam}
For independent exponential variables:
\[
\min(X,Y)\sim \operatorname{Exp}(\lambda_X+\lambda_Y).
\]
Minimum means the first one to happen, so the rates add.

\subsection*{Common trap}
Do not add the expected values of \(X\) and \(Y\). The minimum is not \(X+Y\). For the minimum, add the rates, not the means.

\subsection*{Repeated exam pattern}
No earlier matching pattern found yet in the processed folders. Track this as: minimum of independent exponential variables.

\section*{Uppgift 4}

\subsection*{What this question is really asking}
This question has two separate parts.

Part a is about a discrete random variable. Discrete means it can only take a few separate values:
\[
0,\quad 50,\quad 100.
\]
We need its distribution, expected value, and standard deviation.

Part b is about a taxi cost built from two normally distributed variables:
\[
X=\text{distance in km},\qquad Y=\text{time in minutes}.
\]
We must find the probability that the total cost is between 150 kr and 200 kr.

\subsection*{How to recognize this type of question}
Part a clues:
\begin{itemize}[leftmargin=*]
    \item The possible outcomes are fixed amounts of money.
    \item The text gives percentages and ratio information.
    \item It asks for expected value and standard deviation.
\end{itemize}

Part b clues:
\begin{itemize}[leftmargin=*]
    \item The question says \(X\in N(3,1)\) and \(Y\in N(8,2)\).
    \item That means normal distribution.
    \item The cost is a linear expression: start fee plus km fee plus minute fee.
\end{itemize}

\subsection*{Tiny theory needed}
For a discrete variable:
\[
E(X)=\sum xP(X=x).
\]
This means:

multiply each value by its probability, then add.

For variance:
\[
V(X)=E(X^2)-[E(X)]^2.
\]
Then:
\[
D(X)=\sqrt{V(X)}.
\]

For a linear combination of independent normal variables, the result is also normal.

If:
\[
C=57+14X+9Y,
\]
then:
\[
E(C)=57+14E(X)+9E(Y).
\]
For variance:
\[
V(C)=14^2V(X)+9^2V(Y).
\]
We square the coefficients in variance.

\subsection*{Formula and what everything means}
Discrete expected value:
\[
E(X)=\sum xP(X=x).
\]
Here:
\begin{itemize}[leftmargin=*]
    \item \(x\) is a possible value.
    \item \(P(X=x)\) is the probability of that value.
    \item The sum means add all value-times-probability terms.
\end{itemize}

Normal standardization:
\[
z=\frac{x-\mu}{\sigma}.
\]
Here:
\begin{itemize}[leftmargin=*]
    \item \(x\) is the boundary value we are checking.
    \item \(\mu\) is the mean/center.
    \item \(\sigma\) is the standard deviation/spread.
    \item \(z\) tells how many standard deviations away from the mean we are.
\end{itemize}

\subsection*{Step-by-step solution}
\textbf{Part a: Distribution, expected value, and standard deviation.}

The contribution can be:
\[
0,\quad 50,\quad 100.
\]

The question says \(46\%\) gave nothing:
\[
P(X=0)=0.46.
\]

So the remaining probability is:
\begin{align*}
1-0.46&=0.54.
\end{align*}
This \(0.54\) is split between 50 kr and 100 kr.

The share giving 50 kr is twice as large as the share giving 100 kr. Let:
\[
P(X=100)=q.
\]
Then:
\[
P(X=50)=2q.
\]
Together they must add to \(0.54\):
\begin{align*}
q+2q&=0.54\\
3q&=0.54\\
q&=0.18.
\end{align*}
So:
\[
P(X=100)=0.18,
\]
and:
\[
P(X=50)=2q=2\cdot 0.18=0.36.
\]

The distribution is:
\[
\begin{array}{c|ccc}
x & 0 & 50 & 100\\ \hline
P(X=x) & 0.46 & 0.36 & 0.18
\end{array}
\]

Now expected value:
\begin{align*}
E(X)
&=0\cdot 0.46+50\cdot 0.36+100\cdot 0.18\\
&=0+18+18\\
&=36.
\end{align*}
So the average contribution is 36 kr.

Now calculate \(E(X^2)\):
\begin{align*}
E(X^2)
&=0^2\cdot 0.46+50^2\cdot 0.36+100^2\cdot 0.18\\
&=0+2500\cdot 0.36+10000\cdot 0.18\\
&=0+900+1800\\
&=2700.
\end{align*}

Variance:
\begin{align*}
V(X)
&=E(X^2)-[E(X)]^2\\
&=2700-36^2\\
&=2700-1296\\
&=1404.
\end{align*}

Standard deviation:
\begin{align*}
D(X)
&=\sqrt{1404}\\
&\approx 37.47.
\end{align*}

\textbf{Part b: Taxi cost probability.}

The cost has:
\begin{itemize}[leftmargin=*]
    \item 57 kr start fee,
    \item 14 kr per km,
    \item 9 kr per minute.
\end{itemize}

So the total cost is:
\[
C=57+14X+9Y.
\]

We use the standard convention \(N(\mu,\sigma)\), so:
\[
X\sim N(3,1),\qquad Y\sim N(8,2).
\]
This means:
\[
E(X)=3,\quad D(X)=1,
\]
and:
\[
E(Y)=8,\quad D(Y)=2.
\]

Find the mean of \(C\):
\begin{align*}
E(C)
&=57+14E(X)+9E(Y)\\
&=57+14\cdot 3+9\cdot 8\\
&=57+42+72\\
&=171.
\end{align*}

Now find the variance of \(C\). We assume \(X\) and \(Y\) are independent:
\begin{align*}
V(C)
&=14^2V(X)+9^2V(Y).
\end{align*}
Since standard deviation of \(X\) is \(1\), variance is \(1^2=1\). Since standard deviation of \(Y\) is \(2\), variance is \(2^2=4\):
\begin{align*}
V(C)
&=14^2\cdot 1+9^2\cdot 4\\
&=196+324\\
&=520.
\end{align*}
So:
\[
D(C)=\sqrt{520}\approx 22.80.
\]

Therefore:
\[
C\sim N(171,\sqrt{520}).
\]

We want:
\[
P(150\le C\le 200).
\]
Standardize both endpoints:
\begin{align*}
z_1
&=\frac{150-171}{\sqrt{520}}\\
&\approx -0.921,
\end{align*}
and
\begin{align*}
z_2
&=\frac{200-171}{\sqrt{520}}\\
&\approx 1.272.
\end{align*}

So:
\begin{align*}
P(150\le C\le 200)
&=P(-0.921\le Z\le 1.272)\\
&=\Phi(1.272)-\Phi(-0.921)\\
&\approx 0.898-0.179\\
&\approx 0.720.
\end{align*}

\subsection*{Answer}
\[
\boxed{
\begin{array}{c|ccc}
x & 0 & 50 & 100\\ \hline
P(X=x) & 0.46 & 0.36 & 0.18
\end{array}}
\]
\[
\boxed{E(X)=36\text{ kr}},\qquad
\boxed{D(X)\approx 37.47\text{ kr}}
\]
\[
\boxed{P(150\le C\le 200)\approx 0.720}
\]

\subsection*{What to memorize for the exam}
For discrete variables:
\[
E(X)=\sum xP(X=x),\qquad V(X)=E(X^2)-[E(X)]^2.
\]
For linear normal combinations:
\[
\text{means add directly, variances add with squared coefficients.}
\]

\subsection*{Common trap}
Do not add standard deviations directly. In part b, \(14X+9Y\) means the variance uses \(14^2\) and \(9^2\), not \(14\) and \(9\).

\subsection*{Repeated exam pattern}
Part a is a discrete-distribution-from-word-proportions pattern. Part b is the same linear-combination idea as 2025-08, Uppgift 3c, and the normal linear combination pattern tracked from this exam.

\section*{Uppgift 5}

\subsection*{What this question is really asking}
This is a maximum likelihood question, often shortened to ML.

The random variable \(X\) counts how many successful free throws happen before the first missed free throw.

The probability function is:
\[
p_X(k)=p^k(1-p),\qquad k=0,1,2,\ldots
\]
The unknown number is \(p\), the probability of making one free throw.

The question asks:
\begin{itemize}[leftmargin=*]
    \item Derive the ML estimate of \(p\).
    \item Use the observed data to calculate the estimate.
\end{itemize}

\subsection*{How to recognize this type of question}
The clues are:
\begin{itemize}[leftmargin=*]
    \item It says ``ML-skattning'', meaning maximum likelihood estimate.
    \item A probability function with an unknown parameter \(p\) is given.
    \item A sample \(x_1,\ldots,x_n\) is given.
\end{itemize}

So we use the ML routine:
\[
\text{likelihood} \rightarrow \text{log-likelihood} \rightarrow \text{differentiate} \rightarrow \text{solve}.
\]

\subsection*{Tiny theory needed}
The likelihood is the probability of seeing the data, treated as a function of the unknown parameter \(p\).

Because the observations are independent, we multiply their probabilities:
\[
L(p)=\prod_{i=1}^n p^{x_i}(1-p).
\]

We take logs because products are hard to differentiate, but logs turn products into sums.

\subsection*{Formula and what everything means}
Probability function:
\[
p_X(k)=p^k(1-p).
\]
Here:
\begin{itemize}[leftmargin=*]
    \item \(k\) is the number of successes before the first miss.
    \item \(p\) is the probability of success on one shot.
    \item \(1-p\) is the probability of the first miss.
\end{itemize}

Likelihood:
\[
L(p)=\prod_{i=1}^n p^{x_i}(1-p).
\]
Here:
\begin{itemize}[leftmargin=*]
    \item \(x_i\) is one observed result.
    \item \(n\) is the number of observations.
    \item \(\prod\) means multiply all the terms.
\end{itemize}

\subsection*{Step-by-step solution}
\textbf{Part a: Derive the ML estimate.}

Start with the likelihood:
\begin{align*}
L(p)
&=\prod_{i=1}^n p^{x_i}(1-p).
\end{align*}
All the powers of \(p\) combine:
\[
p^{x_1}p^{x_2}\cdots p^{x_n}=p^{x_1+x_2+\cdots+x_n}.
\]
Let:
\[
S=\sum_{i=1}^n x_i.
\]
This means \(S\) is the sum of all observed values.

There is one \((1-p)\) factor for each observation, so there are \(n\) of them:
\[
(1-p)^n.
\]
Therefore:
\[
L(p)=p^S(1-p)^n.
\]

Now take the natural logarithm:
\begin{align*}
\ell(p)
&=\ln L(p)\\
&=\ln\left(p^S(1-p)^n\right).
\end{align*}
Use log rules:
\begin{align*}
\ell(p)
&=S\ln p+n\ln(1-p).
\end{align*}

Now differentiate:
\[
\frac{d}{dp}(S\ln p)=\frac{S}{p}.
\]
Also:
\[
\frac{d}{dp}\bigl(n\ln(1-p)\bigr)=-\frac{n}{1-p}.
\]
So:
\[
\ell'(p)=\frac{S}{p}-\frac{n}{1-p}.
\]

Set the derivative equal to zero:
\begin{align*}
\frac{S}{p}-\frac{n}{1-p}&=0.
\end{align*}
Move the second term to the other side:
\begin{align*}
\frac{S}{p}&=\frac{n}{1-p}.
\end{align*}
Cross-multiply:
\begin{align*}
S(1-p)&=np.
\end{align*}
Expand the left side:
\begin{align*}
S-Sp&=np.
\end{align*}
Move the \(p\)-terms to the right:
\begin{align*}
S&=np+Sp\\
S&=p(n+S).
\end{align*}
Solve for \(p\):
\begin{align*}
\hat p_{ML}
&=\frac{S}{S+n}.
\end{align*}
Since \(S=\sum x_i\), we can write:
\[
\hat p_{ML}=\frac{\sum_{i=1}^n x_i}{\sum_{i=1}^n x_i+n}.
\]

\textbf{Part b: Insert the data.}

The observed data are:
\[
2,\;0,\;9,\;1,\;6,\;7,\;5,\;3.
\]
There are:
\[
n=8
\]
observations.

Add the observations:
\begin{align*}
S
&=2+0+9+1+6+7+5+3\\
&=33.
\end{align*}

Now insert \(S=33\) and \(n=8\):
\begin{align*}
\hat p_{ML}
&=\frac{33}{33+8}\\
&=\frac{33}{41}\\
&\approx 0.805.
\end{align*}

\subsection*{Answer}
\[
\boxed{\hat p_{ML}=\frac{\sum_{i=1}^n x_i}{\sum_{i=1}^n x_i+n}}
\]
\[
\boxed{\hat p_{ML}=\frac{33}{41}\approx 0.805}
\]

\subsection*{What to memorize for the exam}
For ML questions, the routine is:
\[
L \rightarrow \ln L \rightarrow \text{differentiate} \rightarrow \text{set equal to zero} \rightarrow \text{solve}.
\]
For this specific distribution:
\[
\hat p=\frac{\sum x_i}{\sum x_i+n}.
\]

\subsection*{Common trap}
Do not stop after writing the likelihood. The estimate comes from differentiating the log-likelihood and solving.

\subsection*{Repeated exam pattern}
This is the same broad ML pattern as 2025-08, Uppgift 4. The distribution is different, but the workflow is the same.

\section*{Uppgift 6}

\subsection*{What this question is really asking}
This question is about a sample of daily parking income.

We are given:
\[
n=30,\qquad \bar{x}=1260,\qquad s=150.
\]
Here:
\begin{itemize}[leftmargin=*]
    \item \(n=30\) means 30 weekdays were observed.
    \item \(\bar{x}=1260\) means the sample average was 1260 kr.
    \item \(s=150\) means the sample standard deviation was 150 kr.
\end{itemize}

The question asks for:
\begin{itemize}[leftmargin=*]
    \item a 90\% confidence interval for the true mean daily income,
    \item hypotheses for testing the consultant's claim \(\mu=1280\),
    \item a hypothesis test at significance level 10\%.
\end{itemize}

\subsection*{How to recognize this type of question}
The clues are:
\begin{itemize}[leftmargin=*]
    \item It says ``konfidensintervall'', so part a is a confidence interval.
    \item It gives a sample mean and sample standard deviation.
    \item It says the data are from a normal distribution.
    \item It asks whether a claim about the mean can be rejected, so parts b and c are a hypothesis test.
\end{itemize}

Because the population standard deviation \(\sigma\) is not given, we use the \(t\)-distribution.

\subsection*{Tiny theory needed}
A confidence interval gives a range of plausible values for the true mean \(\mu\).

Here we do not know the true standard deviation \(\sigma\). We only know the sample standard deviation \(s=150\). That is why we use a \(t\)-value instead of a normal \(z\)-value.

A hypothesis test compares:
\begin{itemize}[leftmargin=*]
    \item \(H_0\): the default claim,
    \item \(H_1\): the alternative.
\end{itemize}

The consultant claimed:
\[
\mu=1280.
\]
The question asks whether this claim is reasonable or can be rejected. Since it does not say specifically ``larger'' or ``smaller'', we use a two-sided test:
\[
H_1:\mu\ne 1280.
\]

\subsection*{Formula and what everything means}
Confidence interval for mean with unknown \(\sigma\):
\[
\bar{x}\pm t_{\alpha/2,n-1}\frac{s}{\sqrt{n}}.
\]
Here:
\begin{itemize}[leftmargin=*]
    \item \(\bar{x}\) is the sample mean.
    \item \(t_{\alpha/2,n-1}\) is the \(t\)-critical value.
    \item \(n-1\) is degrees of freedom.
    \item \(s\) is the sample standard deviation.
    \item \(s/\sqrt{n}\) is the standard error, the spread of the sample mean.
\end{itemize}

Test statistic:
\[
T=\frac{\bar{x}-\mu_0}{s/\sqrt{n}}.
\]
Here:
\begin{itemize}[leftmargin=*]
    \item \(\mu_0\) is the value claimed in \(H_0\), here \(1280\).
    \item \(T\) tells how many standard errors the observed mean is from the claimed mean.
\end{itemize}

\subsection*{Step-by-step solution}
\textbf{Part a: 90\% confidence interval.}

Given:
\[
n=30,\qquad \bar{x}=1260,\qquad s=150.
\]

A 90\% interval means:
\[
\alpha=0.10.
\]
Because the interval has two sides, split this into two tails:
\[
\alpha/2=0.05.
\]

Degrees of freedom:
\[
n-1=30-1=29.
\]

From a \(t\)-table:
\[
t_{0.95,29}\approx 1.699.
\]

Now calculate the standard error:
\begin{align*}
\frac{s}{\sqrt{n}}
&=\frac{150}{\sqrt{30}}\\
&\approx \frac{150}{5.477}\\
&\approx 27.386.
\end{align*}

Now calculate the margin of error:
\begin{align*}
1.699\cdot 27.386
&\approx 46.53.
\end{align*}

So the confidence interval is:
\begin{align*}
1260\pm 46.53.
\end{align*}
Lower endpoint:
\begin{align*}
1260-46.53&=1213.47.
\end{align*}
Upper endpoint:
\begin{align*}
1260+46.53&=1306.53.
\end{align*}
So:
\[
(1213.47,\;1306.53).
\]

\textbf{Part b: State hypotheses.}

The consultant claimed:
\[
\mu=1280.
\]
So the null hypothesis is:
\[
H_0:\mu=1280.
\]

The question asks whether the claim is reasonable or can be rejected. It does not say we only care about higher or lower, so use a two-sided alternative:
\[
H_1:\mu\ne 1280.
\]

\textbf{Part c: Perform the hypothesis test.}

Use:
\[
T=\frac{\bar{x}-\mu_0}{s/\sqrt{n}}.
\]
Insert values:
\begin{align*}
T
&=\frac{1260-1280}{150/\sqrt{30}}\\
&=\frac{-20}{27.386}\\
&\approx -0.730.
\end{align*}

At significance level \(10\%\), with a two-sided test and \(29\) degrees of freedom, the critical values are:
\[
\pm t_{0.95,29}\approx \pm 1.699.
\]

The rejection regions are:
\[
T<-1.699
\]
or
\[
T>1.699.
\]

Our value is:
\[
T=-0.730.
\]
This is between \(-1.699\) and \(1.699\):
\[
-1.699<-0.730<1.699.
\]
So we do not reject \(H_0\).

Plain-language check: the claimed value \(1280\) is inside the 90\% confidence interval:
\[
1213.47<1280<1306.53.
\]
So the consultant's claim is plausible from this sample.

\subsection*{Answer}
\[
\boxed{\text{90\% confidence interval for }\mu:\ (1213.47,\;1306.53)\text{ kr}}
\]
\[
\boxed{H_0:\mu=1280,\qquad H_1:\mu\ne 1280}
\]
\[
\boxed{T\approx -0.730}
\]
\[
\boxed{\text{Do not reject }H_0. \text{ The consultant's claim cannot be rejected.}}
\]

\subsection*{What to memorize for the exam}
If the question gives \(s\), the sample standard deviation, and does not give the true \(\sigma\), use a \(t\)-method:
\[
\bar{x}\pm t\frac{s}{\sqrt{n}}.
\]
For a hypothesis test:
\[
T=\frac{\bar{x}-\mu_0}{s/\sqrt{n}}.
\]

\subsection*{Common trap}
Do not use a \(z\)-test just because the sample size is 30. The safer course-exam rule is: if \(\sigma\) is unknown and you use \(s\), use \(t\).

\subsection*{Repeated exam pattern}
This is related to 2025-08, Uppgift 6 because both are hypothesis tests for a mean. The difference is that 2026-01 uses unknown \(\sigma\) and a two-sided \(t\)-test, while 2025-08 uses known \(\sigma\) and a one-sided \(z\)-test.

\section*{Mini-cheatsheet from this exam}

\subsection*{Probability}
Complement:
\[
P(A^c)=1-P(A).
\]

Conditional probability:
\[
P(B\mid A)=\frac{P(A\cap B)}{P(A)}.
\]

Independence:
\[
P(A\cap B)=P(A)P(B).
\]

\subsection*{Continuous density}
To make a density:
\[
\int f(x)\,dx=1.
\]

CDF:
\[
F_X(x)=P(X\le x)=\int_{-\infty}^{x}f(t)\,dt.
\]

Right-tail probability:
\[
P(X>c)=1-F_X(c).
\]

Variance:
\[
V(X)=E(X^2)-[E(X)]^2.
\]

\subsection*{Exponential distribution}
If \(X\sim \operatorname{Exp}(\lambda)\):
\[
P(X>t)=e^{-\lambda t},
\qquad
E(X)=\frac{1}{\lambda}.
\]

For independent exponentials:
\[
\min(X,Y)\sim \operatorname{Exp}(\lambda_X+\lambda_Y).
\]

\subsection*{Discrete variables}
Expected value:
\[
E(X)=\sum xP(X=x).
\]

Second moment:
\[
E(X^2)=\sum x^2P(X=x).
\]

Variance:
\[
V(X)=E(X^2)-[E(X)]^2.
\]

\subsection*{Linear combinations}
Mean:
\[
E(a+bX+cY)=a+bE(X)+cE(Y).
\]

Variance, if independent:
\[
V(a+bX+cY)=b^2V(X)+c^2V(Y).
\]

\subsection*{Normal standardization}
\[
z=\frac{x-\mu}{\sigma}.
\]

\subsection*{Maximum likelihood}
ML workflow:
\[
L \rightarrow \ln L \rightarrow \text{differentiate} \rightarrow \text{set equal to zero} \rightarrow \text{solve}.
\]

For \(p_X(k)=p^k(1-p)\):
\[
\hat p=\frac{\sum x_i}{\sum x_i+n}.
\]

\subsection*{Confidence interval and test for mean, unknown sigma}
Confidence interval:
\[
\bar{x}\pm t_{\alpha/2,n-1}\frac{s}{\sqrt{n}}.
\]

Test statistic:
\[
T=\frac{\bar{x}-\mu_0}{s/\sqrt{n}}.
\]

\end{document}
